\section*{3.8 OP-closure phase: two-core reduction and closure discipline}
\addcontentsline{toc}{section}{3.8 OP-closure phase: two-core reduction and closure discipline}
The next cycle is now promoted from a planning note to a working mathematical stage. The point of 3.8 is still \emph{not} to lift the deferred QAM-, ZFM-, or language-layer upward. Instead, 3.8 treats the remaining foundational closure burden as a two-core residual program over the already fixed residual-record interface. The guiding question is therefore no longer ``what further visible structure is missing?'', but rather ``which residual OP cores must still be closed before the historical OP layer can be regarded as exhausted at the foundational level?''

\begin{definition}[Primary closure core and secondary closure family]
\label{def:mainline-3800}
In the present stabilized reading, the \emph{primary closure core} consists of exactly two residual tasks:
\begin{enumerate}[label=(\roman*), leftmargin=2em]
  \item the OP~1 prime-lattice / continuum-limit correction problem for the matter-frame value $\bsym\approx 0.5493$;
  \item the OP~5 step~(B) Theta-sum / Theta--$R$ balance problem at the photon-frame cutoff.
\end{enumerate}
The \emph{secondary closure family} consists of the remaining open blocks --- controlled emergence / field-equation derivation, mass hierarchy, three-generation splitting, and compact Dixon reformulation --- insofar as their unresolved content is downstream of the same residual-record semantics and does not reopen the visible primitive basis.
\end{definition}

\begin{theorem}[3.8 two-core closure reduction principle]
\label{thm:mainline-3800}
Assume the 3.7 residual downstream reduction principle. Then the foundational historical OP burden of the Companion is reduced to the primary closure core from Definition~\ref{def:mainline-3800} in the following precise sense.
\begin{enumerate}[label=(\roman*), leftmargin=2em]
  \item No already synchronized historical OP outside OP~1 and OP~5 step~(B) remains a primary foundational blocker of the main line.
  \item Any later successful closure of the primary closure core within the existing residual-record semantics closes the remaining historical-interface burden \emph{as such}: after that point, the secondary closure family may still require substantial mathematics, estimates, or phenomenological extraction, but it can no longer be counted as evidence that the visible interface was incomplete.
  \item Consequently, 3.8 may be organized as an OP-closure phase with a strict dependency order: first the OP~1 correction problem and the OP~5 step~(B) balance problem, then the secondary closure family, and only afterwards any renewed consideration of deferred language growth.
\end{enumerate}
\end{theorem}

\begin{proof}
Item~(i) is exactly the synchronized status reached in Sections~3.5--3.7: OP~2 is closed as a separate historical operator problem, OP~3 is closed at the present structural level, and OP~5 step~(A) is already removed from the live main-line burden. What remains as genuinely primary closure work is therefore the residual OP~1 correction problem together with OP~5 step~(B). Item~(ii) follows from the residual-record exact representability and conservative-frontend discipline: once these two residual cores are closed without introducing any new visible primitive, the historical reason for reopening the visible interface disappears. The remaining tasks of controlled emergence, spectral assignment, generation splitting, and Dixon packaging may still be mathematically deep, but they then belong to downstream completion and extraction over a fixed interface rather than to unresolved foundational interface formation. Item~(iii) is therefore the admissible work-ordering rule for the next stage.
\end{proof}

\begin{corollary}[3.8 closure discipline rule]
\label{cor:mainline-3800}
At the start of the 3.8 cycle, the main line should treat every proposed detour according to one question first: does it close part of the primary closure core? If not, then in the absence of a demonstrated dependency it is secondary to the OP~1 / OP~5 step~(B) program. In particular, deferred QAM-, ZFM-, or typed-language growth remains non-urgent unless one of the two primary closure cores genuinely requires it.
\end{corollary}

\begin{proof}
This is the direct operational reading of Theorem~\ref{thm:mainline-3800}. The theorem does not claim that the secondary closure family is trivial, only that it is no longer first in the dependency order of the foundational closure program.
\end{proof}

\paragraph{3.8.0 working remark.}
The practical effect of this step is to make the next cycle much narrower. The burden of proof is now concentrated on two residual cores instead of a diffuse list of historical open fronts. That does not yet solve OP~1 or OP~5 step~(B), but it does mean that a successful 3.8 line can be judged by a much clearer criterion: did we close one of the two cores, reduce one of them to a still smaller kernel, or prove that some secondary task is in fact dependent on one of them?

\paragraph{3.8.0 status note.}
This is a local working-stage promotion, not yet a public-release claim. What is fixed here is the dependency structure of the next closure cycle: the historical OP layer is now to be attacked through a two-core program, with the remaining open blocks treated as downstream completion unless a contrary dependency is demonstrated.

\paragraph{C.28. Two-core reduction logic.}
Once the residual-record interface is fixed and later frontend growth is constrained to be conservative, a historical-open-problem list can be partially ordered by dependency on the visible interface. The point of the two-core reduction is that the remaining foundational closure burden is no longer distributed across many incomparable items: it is concentrated in the OP~1 correction problem and the OP~5 step~(B) balance problem.

\paragraph{C.29. What remains open after two-core reduction.}
The two-core reduction does \emph{not} claim that field equations, mass hierarchy, generation splitting, or Dixon packaging are solved. It claims something narrower and more structural: those tasks no longer appear first in the dependency ordering of foundational closure, and they do not presently justify reopening the visible primitive basis. They remain mathematically serious, but downstream.

\subsection*{3.8.1 residual-controlled emergence equation}
The two-core reduction makes it possible to revisit the field-equation / controlled-emergence question without reopening the visible interface. The point is not to replace the structural master equation by a second fundamental equation, but to write the physical readout problem in a form whose unresolved content is explicitly concentrated into the same two residual cores that already govern the 3.8 closure program.

\begin{definition}[Residual-controlled physical interface form]
\label{def:mainline-3810}
Let \(\mathcal E_{\mathrm{str}}[\Psi]\) denote the structural master side of the present framework, schematically read as
\[
  \mathcal E_{\mathrm{str}}[\Psi]
  :=
  \mathfrak G[\Psi,g]
  + \mathfrak T_{\mathrm{tri}}[\Psi]
  + \mathfrak H_{\mathrm{HC}}[\Psi]
  + \mathfrak Q_{\mathrm{cl}}[\Psi]
  + \mathfrak B_{\beta}[\Psi].
\]
A \emph{residual-controlled physical interface form} is an admissible readout equation of the shape
\[
  \Pi_{\mathrm{phys}}\bigl(\mathcal E_{\mathrm{str}}[\Psi]\bigr)
  =
  \mathcal E_{1}[\Psi]
  +
  \mathcal E_{5}[\Psi],
\]
where \(\Pi_{\mathrm{phys}}\) is a closure-compatible physical readout / observer projection, \(\mathcal E_{1}\) is a remainder family carrying the OP~1 prime-lattice / continuum-limit correction burden, and \(\mathcal E_{5}\) is a remainder family carrying the OP~5 step~(B) Theta-sum / Theta--\(R\) balance burden.
\end{definition}

\begin{theorem}[3.8.1 residual-controlled emergence principle]
\label{thm:mainline-3810}
Assume the Master probe field equation (Theorem~\texttt{thm:master-field-equation} in the active-core PDF), the residual-record exact representability principle, and the 3.8 two-core closure reduction principle. Then the unresolved field-equation / controlled-emergence burden of the present Companion admits a residual-controlled physical interface form in the sense of Definition~\ref{def:mainline-3810}. More precisely:
\begin{enumerate}[label=(\roman*), leftmargin=2em]
  \item the structural master side may be projected to an admissible physical readout equation without introducing any new visible primitive beyond the residual-record interface;
  \item every unresolved contribution that still survives in the present stabilized reading can be absorbed into the two remainder families \(\mathcal E_{1}\) and \(\mathcal E_{5}\), corresponding respectively to the OP~1 correction problem and the OP~5 step~(B) balance problem;
  \item consequently, the field-equation question is no longer a separate primary foundational front parallel to the OP program, but an OP-first residual task over the already fixed interface.
\end{enumerate}
\end{theorem}

\begin{proof}
Theorem~\texttt{thm:master-field-equation} in the active-core PDF already fixes a closure-compatible structural master equation for the probe/field sector. Sections~3.5 and~3.6 then force all later readout and frontend organization to remain conservative over the residual record, so no admissible physical projection may introduce a genuinely new visible primitive.

Theorem~\ref{thm:mainline-3800} further reduces the remaining foundational burden to two cores only: the OP~1 prime-lattice/continuum-limit correction and the OP~5 step~(B) Theta-balance problem. Therefore any unresolved contribution in the controlled-emergence reading of the structural master equation must, in the present stabilized reading, belong to one of these two residual families. Writing the projected equation in the form
\[
  \Pi_{\mathrm{phys}}\bigl(\mathcal E_{\mathrm{str}}[\Psi]\bigr)
  =
  \mathcal E_{1}[\Psi]+\mathcal E_{5}[\Psi]
\]
is thus not a new ontological postulate but a bookkeeping theorem for the remaining closure burden. Items~(i)--(iii) are exactly the content of this repackaging.
\end{proof}

\begin{corollary}[3.8.1 OP-first field-equation strategy]
\label{cor:mainline-3810}
A later controlled derivation of Einstein-, Dirac-, or QFT-type projection formulae need not be treated as a third independent foundational closure front, provided it remains conservative over the residual record and its unresolved content continues to factor through \(\mathcal E_{1}\) and \(\mathcal E_{5}\). In particular, new language growth such as QAM remains non-urgent for the field-equation question unless one of the two residual families genuinely requires it.
\end{corollary}

\begin{proof}
This is the direct operational reading of Theorem~\ref{thm:mainline-3810}. Once the unresolved physical-interface burden is concentrated in \(\mathcal E_{1}\) and \(\mathcal E_{5}\), the field-equation task can only regain primary-foundational status by exhibiting a new visible primitive or by proving that one of the two remainder families cannot be represented within the existing interface. Neither move is currently licensed.
\end{proof}

\paragraph{3.8.1 working remark.}
This step does not claim a finished Einstein equation, a finished Dirac projection, or a closed emergence derivation with all constants fixed. What it does claim is narrower and more useful for 3.8: the field-equation problem can now be attacked as a residual equation with two named remainder channels instead of as an independent search for a new formal foundation.

\paragraph{3.8.1 status note.}
At the level of project strategy, this is the point where a new field equation becomes helpful and QAM remains secondary. The residual-controlled interface equation is useful because it exports the field-equation burden directly into the two live OP cores. QAM would mainly reorganize notation and proof syntax, whereas the present step reorganizes the remaining mathematics.

\paragraph{C.30. Residual-controlled field-equation bookkeeping.}
The residual-controlled interface form should not be misread as a replacement of the structural master equation by a second foundational law. Its purpose is to expose the unresolved physical readout burden as a sum of two residual families already known from the OP-closure program. This is why it can help 3.8 directly even before any later language engineering is undertaken.

\paragraph{C.31. Why QAM remains deferred here.}
In the present stabilized reading, QAM would primarily compress and reorganize proofs. By contrast, the residual-controlled interface equation narrows the actual unresolved mathematics by identifying where the field-equation burden still lives. As long as that burden is concentrated in \(\mathcal E_{1}\) and \(\mathcal E_{5}\), direct OP closure remains the preferred next move.


\subsection*{3.8.2 OP~1 scalar residual channel}
The next useful closure move is to compress OP~1 to the smallest genuinely live mathematical object. The present document already contains the continuum fixed-point candidate, the self-consistent physical matching system, and the statement that the surviving burden is the prime-lattice / continuum-limit correction. What is added here is the observation that, in the present stabilized reading, this burden can be written as a single scalar residual channel rather than as a diffuse family of unrelated corrections.

\begin{definition}[OP~1 scalar residual profile]
\label{def:mainline-3820}
For \(\mu>1\), define the continuum-induced matter-frame candidate
\[
  \beta_{\mathrm{cont}}(\mu):=\sqrt{\frac{\mu-1}{\mu+1}}
\]
and the spectral correction residual
\[
  \mathcal R_{1}(\mu)
  :=
  \log\!\bigl(\sqrt{3}\,\beta_{\mathrm{cont}}(\mu)\bigr)
  +
  \frac{\gamma_1}{4\pi\mu^5}.
\]
The equation \(\mathcal R_{1}(\mu)=0\) is called the \emph{OP~1 scalar residual channel}. It is the logarithmic form of the self-consistent matching condition
\[
  \beta_{\mathrm{cont}}(\mu)
  =
  \frac{1}{\sqrt{3}}\exp\!\Bigl(-\frac{\gamma_1}{4\pi\mu^5}\Bigr).
\]
\end{definition}

\begin{theorem}[3.8.2 OP~1 single-channel reduction principle]
\label{thm:mainline-3820}
Assume the Tribonacci continuum fixed-point theorem, the self-consistent OP~1 matching system (\texttt{eq:SC1}, active-core PDF)--(\texttt{eq:SC2}, active-core PDF), and the 3.8.1 residual-controlled emergence principle. Then, in the present stabilized reading:
\begin{enumerate}[label=(\roman*), leftmargin=2em]
  \item the remaining OP~1 burden is completely represented by the scalar residual profile \(\mathcal R_{1}(\mu)\) of Definition~\ref{def:mainline-3820};
  \item closing OP~1 no longer requires a new visible primitive, a new field equation, or a new language layer, but only a first-principles control or derivation of the correction law encoded by \(\mathcal R_{1}(\mu)\);
  \item any later closure of OP~1 that is admissible in the current framework must factor through the vanishing or controlled estimation of this scalar residual channel.
\end{enumerate}
\end{theorem}

\begin{proof}
The continuum side of OP~1 is already fixed by Theorem~\texttt{thm:tribonacci} in the active-core PDF: the uncorrected self-referential frame candidate is the Tribonacci value, equivalently read through the continuum expression \(\beta_{\mathrm{cont}}(\mu)=\sqrt{(\mu-1)/(\mu+1)}\). The physical matching side is already fixed by the self-consistent system (\texttt{eq:SC1}, active-core PDF)--(\texttt{eq:SC2}, active-core PDF), whose unresolved content is precisely the prime-lattice / continuum-limit correction law. Taking logarithms rewrites the matching problem as the vanishing of the single scalar expression \(\mathcal R_{1}(\mu)\). This loses no information because \(\mu>1\) and the exponential side of (\texttt{eq:SC1}, active-core PDF) is positive. By the 3.8.1 residual-controlled emergence principle, the field-equation burden cannot add a third independent OP~1 channel; it can only feed into the same residual family \(\mathcal E_{1}\). Hence the entire remaining OP~1 content is exhausted by the scalar residual channel \(\mathcal R_{1}(\mu)\), and any admissible closure must proceed by deriving, controlling, or eliminating that residual.
\end{proof}

\begin{corollary}[3.8.2 OP~1 closure test]
\label{cor:mainline-3820}
Within the present Companion, OP~1 is closed as soon as one of the following is obtained without introducing a new visible primitive basis:
\begin{enumerate}[label=(\alph*), leftmargin=2em]
  \item a first-principles derivation showing \(\mathcal R_{1}(\gamma_L)=0\) exactly;
  \item a structural derivation of the prime-lattice correction law that identifies \(\mathcal R_{1}\) with an already controlled arithmetic or spectral remainder and proves its vanishing;
  \item a proof that the residual channel is bounded in a way that collapses the remaining numerical gap and stabilizes the same \(\beta_{\mathrm{sym}}\)-value uniquely.
\end{enumerate}
In particular, OP~1 is now a one-channel closure problem.
\end{corollary}

\begin{proof}
This is the direct operational reading of Theorem~\ref{thm:mainline-3820}. Once all remaining OP~1 content is represented by a single scalar residual channel, closing that channel by exact vanishing, structural identification, or sufficiently strong bound is enough to close OP~1 within the present stabilized reading.
\end{proof}

\paragraph{3.8.2 working remark.}
This step does not yet prove the prime-lattice correction from first principles. What it does prove is that the surviving OP~1 mathematics has become one-dimensional in a strong structural sense: there is no longer a family of incomparable OP~1 subproblems, but a single residual channel to be killed or identified.

\paragraph{3.8.2 status note.}
This is the first real 3.8 closure compression after the two-core reduction. It is also a useful diagnostic: if OP~1 later resists closure, the obstruction should now show up as a failure to control \(\mathcal R_{1}\), rather than as a vague demand for new visible machinery.

\paragraph{C.32. Why the OP~1 residual channel is scalar.}
The two ingredients entering the current OP~1 matching problem are already fixed: the continuum frame candidate and the prime-lattice spectral correction ansatz. Their discrepancy can therefore be measured by a single logarithmic mismatch. This is why the surviving burden is scalar even though its derivation may still be deep.

\paragraph{C.33. What 3.8.2 does \emph{not} claim.}
Section~3.8.2 does not claim a first-principles proof of the \(\mu^5\) exponent, nor a full arithmetic derivation of the prime-lattice correction. It claims something narrower: that all such work is now forced into a single residual channel and no longer distributed across several independent visible fronts.


\subsection*{3.8.3 OP~1 root isolation and near-closure window}
Once OP~1 has been compressed to the scalar residual channel \(\mathcal R_{1}(\mu)\), the next useful move is not yet to claim exact closure, but to show that the residual channel already has a unique physical root in the relevant regime and that the current \(\gamma_L\)-value lies extremely close to it. This upgrades OP~1 from a diffuse correction problem to a sharply localized near-closure problem.

\begin{definition}[OP~1 physical root window]
\label{def:mainline-3830}
Let
\[
  I_{\mathrm{phys}}:=[1.85,1.90].
\]
Within the present Companion, the interval \(I_{\mathrm{phys}}\) is called the \emph{OP~1 physical root window}. It contains the current matter-frame value \(\gamma_L\approx 1.864\) and the single-channel OP~1 residual profile \(\mathcal R_1(\mu)\) of Definition~\ref{def:mainline-3820}.
\end{definition}

\begin{theorem}[3.8.3 OP~1 root isolation principle]
\label{thm:mainline-3830}
Within the OP~1 scalar residual channel of Definition~\ref{def:mainline-3820}, the following hold on the physical root window \(I_{\mathrm{phys}}=[1.85,1.90]\):
\begin{enumerate}[label=(\roman*), leftmargin=2em]
  \item \(\mathcal R_1\) is strictly increasing on \(I_{\mathrm{phys}}\);
  \item \(\mathcal R_1(1.85)<0\) and \(\mathcal R_1(1.90)>0\);
  \item hence there exists a unique root
  \[
    \mu_\ast\in(1.85,1.90)
  \]
  such that \(\mathcal R_1(\mu_\ast)=0\);
  \item numerically,
  \[
    \mu_\ast \approx 1.86366225,
  \]
  so the current matter-frame value \(\gamma_L\approx 1.864\) differs from the unique OP~1 root only by
  \[
    |\gamma_L-\mu_\ast|\approx 3.38\times 10^{-4},
  \]
  i.e. at the level of about \(1.8\times 10^{-4}\) relative to \(\mu_\ast\).
\end{enumerate}
\end{theorem}

\begin{proof}
From Definition~\ref{def:mainline-3820},
\[
  \mathcal R_1(\mu)=\frac12\log\!\left(\frac{3(\mu-1)}{\mu+1}\right)+\frac{\gamma_1}{4\pi\mu^5}.
\]
Therefore
\[
  \mathcal R_1'(\mu)
  =\frac{1}{\mu^2-1}-\frac{5\gamma_1}{4\pi\mu^6}.
\]
On \(I_{\mathrm{phys}}=[1.85,1.90]\), the first term is decreasing and the second term is increasing in absolute value only mildly; direct evaluation at the right endpoint gives
\[
  \mathcal R_1'(1.90)
  =
  \frac{1}{1.90^2-1}-\frac{5\gamma_1}{4\pi\cdot 1.90^6}
  > 0.265-0.255
  > 0.
\]
Hence \(\mathcal R_1'(\mu)>0\) throughout \(I_{\mathrm{phys}}\), so \(\mathcal R_1\) is strictly increasing there. Direct evaluation gives
\[
  \mathcal R_1(1.85)\approx -3.71\times 10^{-3}<0,
  \qquad
  \mathcal R_1(1.90)\approx 9.70\times 10^{-3}>0.
\]
By continuity and strict monotonicity, there is a unique \(\mu_\ast\in(1.85,1.90)\) with \(\mathcal R_1(\mu_\ast)=0\). A numerical solve of the scalar residual equation yields
\[
  \mu_\ast\approx 1.86366225.
\]
Comparing with the current matter-frame value \(\gamma_L\approx 1.864\) gives
\[
  |\gamma_L-\mu_\ast|\approx 3.38\times 10^{-4},
\]
and the corresponding relative mismatch is about \(1.8\times 10^{-4}\). No new visible primitive enters this argument: it is purely an analysis of the already isolated scalar residual channel.
\end{proof}

\begin{corollary}[3.8.3 OP~1 exact-vs-structural remainder]
\label{cor:mainline-3830}
After Theorem~\ref{thm:mainline-3830}, the remaining OP~1 burden is no longer to locate a candidate value for \(\gamma_L\), but only to explain structurally why the exact root \(\mu_\ast\) of the scalar residual channel must coincide with the framework value read as \(\gamma_L\). Equivalently, the surviving task is an exact prime-lattice / continuum derivation of the residual law, not a search for a second numerical branch.
\end{corollary}

\begin{proof}
Theorem~\ref{thm:mainline-3830} proves that the scalar residual channel already has a unique physical root in the relevant regime and that the current \(\gamma_L\)-value lies extremely close to that root. Therefore the numerical candidate is no longer underdetermined. What remains open is the structural identification of that unique root with the framework quantity \(\gamma_L\), i.e. the first-principles explanation of why the residual channel must vanish exactly there.
\end{proof}

\paragraph{3.8.3 working remark.}
This is the first point in the OP~1 program where the residual channel becomes visibly \emph{near-closed}: not because the prime-lattice correction has been derived, but because the residual equation already has a unique physical root and the framework value currently used for \(\gamma_L\) sits essentially on top of it.

\paragraph{3.8.3 status note.}
Section~3.8.3 still does \emph{not} claim exact first-principles closure of OP~1. What it does claim is narrower and stronger than 3.8.2: the unresolved content is now an exact structural identification problem, not a broad numerical search problem.

\paragraph{C.34. Why root isolation matters.}
A scalar residual channel can still hide several distinct physical branches. The value of Theorem~\ref{thm:mainline-3830} is precisely that this no longer happens in the currently relevant OP~1 window: the residual has a unique root there, so the closure burden becomes conceptual and structural rather than branch-selection driven.

\paragraph{C.35. What 3.8.3 still leaves open.}
The exponent \(\mu^5\), the prime-lattice interpretation of the correction, and the exact identification \(\mu_\ast=\gamma_L\) remain open. The present step only shows that these open points now live over a unique near-closure root rather than over a family of competing candidates.



